\begin{sol}[Carro con cilindro rodante acoplado por un resorte]

\subsubsection*{Pregunta 1: Momento de inercia del cilindro}

Con densidad volumétrica $\rho=\dfrac{m}{\pi r^2 h}$, integrando en coordenadas cilíndricas:
$$I = \int_0^h\!\!\int_0^{2\pi}\!\!\int_0^r s^2\,\rho\,s\,ds\,d\varphi\,dz
= \frac{m}{\pi r^2 h}\cdot\frac{r^4}{4}\cdot 2\pi h = \boxed{\tfrac12 mr^2}.$$

\subsubsection*{Pregunta 2: Ecuaciones de movimiento}

\textbf{Cinemática.} El carro está en $\vec{r}_1=x_1\vec{e}_x$ y el centro del cilindro en
$\vec{r}_2=(x_1+x_2)\vec{e}_x$, luego $\vec{a}_1=\ddot{x}_1\vec{e}_x$ y
$\vec{a}_2=(\ddot{x}_1+\ddot{x}_2)\vec{e}_x$. La rodadura sin resbalar del cilindro sobre el
carro relaciona su rotación con el movimiento relativo: $\dot{x}_2 = r\,\omega$, es decir
$\omega=\dot{x}_2/r$ y $\dot\omega=\ddot{x}_2/r$.

\textbf{Fuerzas.} Sobre el sistema actúan la fuerza externa $f_1$, la del resorte
$F_s=-k(x_2-\ell_0)$, la normal en el contacto y la fuerza tangencial (de roce) $R$ entre carro
y cilindro. Aplicando la primera ley de Euler (traslación) al carro y al cilindro:
$$\text{carro:}\quad f_1 + F_s - R = m\ddot{x}_1,\qquad
\text{cilindro:}\quad R - F_s = m(\ddot{x}_1+\ddot{x}_2).$$
Sumando ambas se elimina $R$ y $F_s$:
$$f_1 = m(2\ddot{x}_1 + \ddot{x}_2). \qquad (\ast)$$

\textbf{Euler sobre el cilindro (respecto a su centro).} Sólo $R$ genera torque:
$$I\dot\omega = -rR \;\Longrightarrow\; \tfrac12 mr^2\,\frac{\ddot{x}_2}{r} = -rR
\;\Longrightarrow\; R = -\tfrac12 m\ddot{x}_2.$$
Sustituyendo $R$ en la ecuación del cilindro, con $F_s=-k(x_2-\ell_0)$:
$$-\tfrac12 m\ddot{x}_2 + k(x_2-\ell_0) = m(\ddot{x}_1+\ddot{x}_2).$$
Combinando esta relación con $(\ast)$ se obtiene finalmente:
$$\boxed{\;\ddot{x}_2 = -\frac{1}{2m}f_1 + \frac{k}{m}(x_2-\ell_0)\;}\qquad
\boxed{\;\ddot{x}_1 = \frac{1}{4m}f_1 - \frac{k}{2m}(x_2-\ell_0)\;}$$

\end{sol}
